
\medskip
\noindent \SGARefTarget{sga7.sga7i.eviii.n7-3-4}{7.3.4}.  \underline{Reduction of obstructions of type (\SGARefLink{sga7.sga7i.eviii.n7-1-2}{7.1.2}) to
those of type (\SGARefLink{sga7.sga7i.eviii.n7-1-1}{7.1.1}).}

Suppose, for definiteness, that there is an integer $n>0$ such that
$n\Phi=0$, and let
\[
 q\in Q(S)
\]
be a section of $Q$.  By an abuse of notation, write
\[
 q_0\in\Psi(k)
\]
for the image in $\Psi(k)$ of the value of $q$ at $s$.  Regarding $E_\eta$
as an extension of $P_\eta\times_\eta Q_\eta$ by
$\underline G_{mQ_\eta}$ over $Q_\eta$, pullback gives an extension
\[
 q_\eta^*(E_\eta),\qquad
 \text{an extension of $P_\eta$ by $\underline G_{m\eta}$},
\]
and hence an obstruction
\begin{equation}\tag{\SGARefTarget{sga7.sga7i.eviii.n7-3-4-1}{7.3.4.1}}
 d_q\overset{\mathrm{dfn}}{=}
 d\bigl(q_\eta^*(E_\eta)\bigr):
 \Phi\longrightarrow(\mathbb Q/\mathbb Z)_k.
\end{equation}
With $d=d(E_\eta)$, one has the relation
\begin{equation}\tag{\SGARefTarget{sga7.sga7i.eviii.n7-3-4-2}{7.3.4.2}}
 d_q(x)=d(x,q_0),
\end{equation}
valid for every point $x$ of $\Phi$ with values in a $k$-scheme $T$.

\smallskip
\noindent \SGARefTarget{sga7.sga7i.eviii.n7-3-4-3}{7.3.4.3}.  For every $q\in Q(S)$, this formula determines $d_q$
in terms of the pairing $d$; conversely, it plainly recovers $d$ from the
invariants $d_q$, at least when every point of $\Psi$ comes from a section
of $Q$.  By Hensel's lemma, this condition holds automatically when $S$ is
strictly local.  The pairing $d$ is known once it is known over the
separable closure of $k$, and its formation commutes with etale localization
on $S$ by \SGARefLink{sga7.sga7i.eviii.n7-3-5-3}{7.3.5.3} below.  Thus, after replacing $S$ by its strict
localization if necessary, the determination of $d$ reduces to that of the
invariants $d_q$.

In particular, when every point of $\Psi$ comes from a section $q$ of $Q$,
the pairing $d$ is zero---that is, $E_\eta$ extends to a biextension $E$ of
$(P,Q)$ by $\underline G_m$---if and only if $d_q$ is zero for every
section $q$ of $Q$, that is, if and only if $q_\eta^*(E_\eta)$ extends to
an extension of $P$ by $\underline G_m$.  With no hypothesis on $\Psi$,
the analogous criterion remains valid provided that the $q$ are allowed to
be points with values in local etale $S$-schemes $S'$.

\smallskip
\noindent \SGARefTarget{sga7.sga7i.eviii.n7-3-4-4}{7.3.4.4}.  Conversely, starting from the situation of \SGARefLink{sga7.sga7i.eviii.n7-1}{7.1}(a), the
obstruction (\SGARefLink{sga7.sga7i.eviii.n7-1-1}{7.1.1}) may be interpreted as an obstruction of type (\SGARefLink{sga7.sga7i.eviii.n7-1-2}{7.1.2}):
take $Q=\mathbb Z_S$ and the biextension of $(P_\eta,\mathbb Z_\eta)$
deduced from $E_\eta$.  Indeed, it suffices to apply the preceding
construction with $q$ equal to the section $1$ of $\mathbb Z_S$.

\medskip
\noindent \SGARefTarget{sga7.sga7i.eviii.n7-3-5}{7.3.5}.  \underline{Effect of a base change $S'\to S$.}  Suppose
given a dominant local morphism of traits
\[
 f:S'\longrightarrow S,
\]
and denote by $\eta'$ the generic point of $S'$, by $s'$ its closed point,
and put $k'=k(s')$.  If $S=\operatorname{Spec}(V)$ and
$S'=\operatorname{Spec}(V')$, then $V'$ is a discrete valuation ring
dominating the discrete valuation ring $V$.  Put
\begin{equation}\tag{\SGARefTarget{sga7.sga7i.eviii.n7-3-5-1}{7.3.5.1}}
 e(S'/S)=e(V'/V)=\operatorname{length}_{V'}(V'/\mathfrak m V'),
\end{equation}
where $\mathfrak m$ is the maximal ideal of $V$.

Let now $P$ be as in \SGARefLink{sga7.sga7i.eviii.n7-1}{7.1}(a), let $E_\eta$ be an extension of $P_\eta$ by
$\underline G_{m\eta}$, put $P'=P\times_S S'$, and let $E'_{\eta'}$ be
the extension of $P'_{\eta'}$ by $\underline G_{m\eta'}$ obtained by
pulling back $E_\eta$.  With the notation of \SGARefLink{sga7.sga7i.eviii.n4-1}{4.1}(a), there is a canonical
isomorphism
\[
 \Phi'\overset{\mathrm{dfn}}{=}P'_0/{P'_0}^{\circ}
 \simeq\Phi\otimes_k k'\simeq\Phi_{k'},
\]
and hence one may consider
\[
 d(E_\eta)_{k'}:\Phi'\longrightarrow(\mathbb Q/\mathbb Z)_{k'}.
\]
Then
\begin{equation}\tag{\SGARefTarget{sga7.sga7i.eviii.n7-3-5-2}{7.3.5.2}}
 d(E'_{\eta'})=e(S'/S)\,d(E_\eta)_{k'}.
\end{equation}

Likewise, under the hypotheses of \SGARefLink{sga7.sga7i.eviii.n7-1}{7.1}(b) and with the evident notation,
(\SGARefLink{sga7.sga7i.eviii.n7-3-5-2}{7.3.5.2}) holds with $d(E_\eta)$ now denoting the pairing (\SGARefLink{sga7.sga7i.eviii.n7-1-2}{7.1.2}),
$d(E_\eta)_{k'}$ the pairing
\[
 \Phi'\otimes\Psi'\longrightarrow(\mathbb Q/\mathbb Z)_{k'}
\]
obtained from it by base change $k\to k'$, and $E'_{\eta'}$ the
biextension of $(P'_{\eta'},Q'_{\eta'})$ by
$\underline G_{m\eta'}$ obtained from the biextension $E_\eta$ of
$(P_\eta,Q_\eta)$ by $\underline G_{m\eta}$ through the base change
$\eta'\to\eta$.

\smallskip
\noindent \SGARefTarget{sga7.sga7i.eviii.n7-3-5-3}{7.3.5.3}.  An interesting case is $e(S'/S)=1$, that is,
\[
 \mathfrak m'=\mathfrak m V',
\]
where $\mathfrak m'$ denotes the maximal ideal of $V'$.  This holds in
particular if $V'$ is etale over $V$, or if $V'$ is the henselization,
strict henselization, or completion of $V$.  In this case (\SGARefLink{sga7.sga7i.eviii.n7-3-5-2}{7.3.5.2}) simply
says that \underline{the formation of obstruction (\SGARefLink{sga7.sga7i.eviii.n7-1-1}{7.1.1}), respectively
(\SGARefLink{sga7.sga7i.eviii.n7-1-2}{7.1.2}), commutes with the base change under consideration $S'\to S$}.

\smallskip
\noindent \SGARefTarget{sga7.sga7i.eviii.n7-3-5-4}{7.3.5.4}.  On the other hand, (\SGARefLink{sga7.sga7i.eviii.n7-3-5-2}{7.3.5.2}) in the general case shows
that, for every extension $E_\eta$ as in \SGARefLink{sga7.sga7i.eviii.n7-1}{7.1}(a) (respectively every
biextension $E_\eta$ as in \SGARefLink{sga7.sga7i.eviii.n7-1}{7.1}(b)), there is a \underline{finite} morphism
of traits $S'\to S$ such that $k(\eta')$ is a separable extension of
$k(\eta)$ and such that the extension (respectively biextension)
$E'_{\eta'}$ obtained from $E_\eta$ by base change $\eta'\to\eta$ extends
to an extension (respectively biextension) $E'$ over all of $S'$.  Indeed,
it is enough to choose $S'/S$ so that $e(S'/S)$ is a multiple of the
integer $n>0$ occurring in \SGARefLink{sga7.sga7i.eviii.n7-1}{7.1}, and such an extension can always be found
with $k(\eta')$ separable over $k(\eta)$.
